\documentclass{article}


\usepackage{arxiv}

\usepackage[utf8]{inputenc} % allow utf-8 input
\usepackage[T1]{fontenc}    % use 8-bit T1 fonts
\usepackage{hyperref}       % hyperlinks
\usepackage{url}            % simple URL typesetting
\usepackage{booktabs}       % professional-quality tables
\usepackage{amsfonts}       % blackboard math symbols
\usepackage{amsmath}
\usepackage{nicefrac}       % compact symbols for 1/2, etc.
\usepackage{microtype}      % microtypography
\usepackage{graphicx}
\usepackage{multirow}
\usepackage{float}          % [H] pins tables and figures exactly where they appear
\usepackage{pdflscape}      % landscape pages for the wide appendix tables
\usepackage{longtable}      % appendix tables that span several pages
\graphicspath{ {./images/} }


\title{A Systematic Study of Design Choices for Multi-Label Chest X-ray
Classification on CheXpert}


\author{
 Ma'moun Yosef \\
  Department of Artificial Intelligence\\
  University of Jordan\\
  Amman, Jordan \\
  \texttt{mam0233493@ju.edu.jo} \\
}

\begin{document}
\maketitle

\begin{abstract}
  Results on CheXpert are hard to compare because published setups differ in backbone,
  resolution, uncertain label handling, objective and ensembling all at once. This report
  fixes the pipeline and varies one factor at a time across roughly fifty runs that share
  the same splits, preprocessing and evaluation. Repeating configurations across seeds
  gives a noise floor of 0.001 mean AUROC on a large automatically labelled validation
  split and 0.004 to 0.017 on the two official radiologist sets, which is wider than many
  of the differences the study set out to measure. Read against that floor, input
  resolution up to $768\times640$ and pretraining on chest radiographs are the choices
  that reliably pay, and the architecture matters more than its size. The ImageNet label
  set, label smoothing and direct AUC optimization all fall inside the noise, while the
  handling of uncertain labels is worth up to 0.023 on one backbone. Ensembling pays when
  its members differ: seven runs of one backbone average to less than three different
  backbones do. The best single model reaches 0.9113 mean AUROC on the official test set
  and a plain average of three backbones reaches 0.9174, a paired bootstrap gain of
  $+0.0061$ with a 95\% interval of $[+0.0005, +0.0118]$. The trained models and their
  configurations are released.
\end{abstract}


% keywords can be removed
%\keywords{chest radiography \and CheXpert \and multi-label classification \and transfer learning \and model selection}


\section{Introduction}

CheXpert \cite{irvin2019chexpert} is one of the standard benchmarks for multi-label
chest radiograph classification, and results on it are difficult to compare. Published
setups differ in backbone, input resolution, uncertain label handling, augmentation,
objective and ensembling all at once, so a reported gain cannot be attributed to any
one of those choices.

This report fixes the pipeline and varies one factor at a time. Roughly fifty runs share
the same splits, preprocessing, augmentation and evaluation, so every difference between
them comes from the design choice under test. Repeated seeds give a noise floor to read
those differences against, which turns out to matter more than any single result here.

The contributions are a controlled comparison of the main design choices on CheXpert, a
measurement of how much of the resulting spread is run to run noise, a final ensemble
reaching 0.9174 mean AUROC on the official test set, and a public release of the trained
models and their configurations.

\section{Methods}

\subsection{Dataset and splits}

All experiments use mainly CheXpert \cite{irvin2019chexpert}, restricted to frontal views.
The targets are the five competition pathologies: Atelectasis, Cardiomegaly,
Consolidation, Edema and Pleural Effusion.

The official training was therefore divided into a training part
and a validation part. The division is 90/10 and split by patient, so no patient
appears in both parts. The three official splits are already patient disjoint.
Sizes are given in Table \ref{tab:splits}.

The official validation set holds only 200 studies, which is too small to give a
stable signal, so the larger val19k split was carved out of the training data for
that purpose. Later runs also scored val200 at every validation and used it to pick
the best checkpoint. The official test500 set was never seen during training and is
kept for final reporting only.

These exact splits are shared by every run, so any difference between runs comes from
the design choice and not from the data. A big portion of the runs were trained at $384\times320$
from the downscaled CheXpert release, and some trained above that resolution from
the original full resolution images.

\begin{table}[H]
  \caption{Splits used in every experiment. Counts are frontal images.}
  \centering
  \begin{tabular}{llrl}
    \toprule
    Split    & Source                           & Images  & Labels                \\
    \midrule
    train    & official train, 90\%             & 171,927 & automatic labeler     \\
    val19k   & official train, 10\%             & 19,100  & automatic labeler     \\
    valid200 & official validation, 200 studies & 202     & radiologist consensus \\
    test500  & official test, 500 studies       & 518     & radiologist consensus \\
    \bottomrule
  \end{tabular}
  \label{tab:splits}
\end{table}


\subsection{Uncertain labels}
\label{sec:uncertain}

CheXpert marks each pathology per study as positive, negative, uncertain or blank.
Blank cells are treated as negative in every run. The uncertain cells are the open
question, and five policies were compared:

\begin{itemize}
  \item \textbf{U-Ones}: uncertain becomes positive.
  \item \textbf{U-Zeros}: uncertain becomes negative.
  \item \textbf{Certain only}: uncertain cells are masked out of the loss.
  \item \textbf{U-MultiClass}: the task gets a three way head over negative,
        positive and uncertain.
  \item \textbf{U-SelfTrained}: uncertain training cells take a soft probability
        predicted by a model that was itself trained on certain cells only.
\end{itemize}

Uncertain labels in val19k are always U-Ones for consistency.
The policies are compared in Section \ref{sec:results-uncertain}.

\subsection{Preprocessing and augmentation}

Each image is resized to fit inside the target width and height with its aspect ratio
preserved, then zero padded on the short side to reach the exact size. Nothing is
stretched and nothing is cropped, so the anatomy keeps its true proportions.
Resampling uses area interpolation. Images are normalized with the statistics of the
data each backbone was pretrained on. The base geometry is $384\times320$, and the higher resolution runs
scale it up to $1600\times1312$.

Contrast limited adaptive histogram equalization is a common preprocessing step for
radiographs, so it was included as an ablation with a clip limit of 2.0 and an
$8\times8$ tile grid, and it did not show a meaningful difference on ResNet-50. It is
off everywhere else.

Augmentation is applied to the training split only, never to any validation or test
split. Each transform is drawn independently with its own probability, so a given
image may receive none, one or several of them. The settings are listed in
Table \ref{tab:augmentation} and are the same in every run. Horizontal flips are
excluded, because mirroring a radiograph moves the heart to the wrong side.

\begin{table}[H]
  \caption{Training augmentation. It is applied to the training split only.}
  \centering
  \begin{tabular}{llc}
    \toprule
    Transform       & Range                      & Probability \\
    \midrule
    Rotation        & $\pm 5$ degrees            & 0.2         \\
    Translation     & $\pm 5\%$ of width, height & 0.2         \\
    Scale           & 0.95 to 1.05               & 0.2         \\
    Brightness      & $\pm 10\%$                 & 0.2         \\
    Contrast        & $\pm 10\%$                 & 0.2         \\
    Horizontal flip & not used                   &             \\
    \bottomrule
  \end{tabular}
  \label{tab:augmentation}
\end{table}

\subsection{Backbones and pretraining}

Twelve backbones were fine-tuned, listed in Table \ref{tab:models}. Parameter counts
are the real counts of the built model, head included. All weights come from torchvision or timm
\cite{wightman2019timm} except the two chest X-ray models. Every backbone is fine
tuned end to end, with no frozen layers.

\begin{table}[H]
  \caption{Backbones, parameter counts and the resolutions each was trained at.}
  \centering
  \begin{tabular}{llrl}
    \toprule
    Backbone                                      & Pretraining          & Params (M) & Resolutions                                      \\
    \midrule
    ResNet-50 \cite{he2016resnet}                 & ImageNet-1k          & 23.5       & $384\times320$                                   \\
    DenseNet-121 \cite{huang2017densenet}         & ImageNet-1k          & 7.0        & $384\times320$                                   \\
    DenseNet-201 \cite{huang2017densenet}         & ImageNet-1k          & 18.1       & $384\times320$                                   \\
    ConvNeXt-T \cite{liu2022convnext}             & ImageNet-1k, 22k     & 27.8       & $384\times320$                                   \\
    ConvNeXt-S \cite{liu2022convnext}             & ImageNet-1k          & 49.5       & $384\times320$                                   \\
    ConvNeXt-B \cite{liu2022convnext}             & ImageNet-22k         & 87.6       & $384\times320$, $768\times640$, $1600\times1312$ \\
    ConvNeXt-L \cite{liu2022convnext}             & ImageNet-22k         & 196.2      & $384\times320$, $768\times640$                   \\
    Swin-B \cite{liu2021swin}                     & ImageNet-1k, 22k     & 86.9       & $384\times384$                                   \\
    EfficientNetV2-M \cite{tan2021efficientnetv2} & ImageNet-21k         & 52.9       & $384\times320$                                   \\
    CoAtNet-1 \cite{dai2021coatnet}               & ImageNet-12k         & 41.0       & $384\times320$                                   \\
    ViT-B, Medical-MAE \cite{xiao2023medicalmae}  & chest X-rays, MAE    & 86.0       & $384\times320$ to $768\times640$                 \\
    ViT-B, RAD-DINO \cite{perezgarcia2024raddino} & chest X-rays, DINOv2 & 86.6       & $784\times644$, $1064\times896$                  \\
    \bottomrule
  \end{tabular}
  \label{tab:models}
\end{table}

The ImageNet-22k, 21k and 12k entries were pretrained on the larger label set and
then fine tuned on ImageNet-1k before being used here.

The last two rows are the only backbones that never saw natural images. Medical-MAE
is a ViT-B trained with masked autoencoding \cite{he2022mae} on roughly 0.5 million
chest radiographs, and three of its released checkpoints were compared: the raw
self-supervised backbone, the same backbone after fine tuning on ChestX-ray14
\cite{wang2017chestxray14}, and the one already fine tuned on CheXpert. RAD-DINO
\cite{perezgarcia2024raddino} is a ViT-B trained with DINOv2 \cite{oquab2024dinov2}
on chest radiographs, and it uses its own normalization statistics rather than the
ImageNet ones. A separate run also pretrained ConvNeXt on ChestX-ray14 before moving
to CheXpert, which tests the same idea with a convolutional backbone.

\subsection{Training setup}

Every run uses AdamW, a batch size of 64, mixed precision, and a cosine schedule with
a one epoch linear warmup. The
loss is binary cross entropy on the logits, computed over the five tasks at once and
masked wherever a target is undefined. Training stops early when the monitored mean
AUROC has not improved for a fixed number of validations.

The weight decay is $1\times10^{-5}$ throughout. The runs that start from a backbone
already pretrained on chest X-rays use a lower peak learning rate than the rest,
because those weights already sit close to a good solution and a larger rate would
only undo them. Those runs also use layer-wise learning rate decay of 0.75, so the
rate is multiplied by 0.75 for each block further from the head, together with gradient
clipping at a gradient norm of 1.0. No layer is frozen anywhere. The peak and minimum learning
rate of every run is listed in
Table \ref{tab:allruns}.

Later runs add label smoothing of 0.05 to the targets, which keeps the model from
driving its logits to saturation on labels that are themselves uncertain. An ablation
was done with and without label smoothing, and label smoothing showed a very small
gain, small enough to sit inside the seed noise reported in
Section \ref{sec:results-seed}. It is kept in every later run.

The high resolution vision transformers are trained in two stages. Their position
embeddings are tied to the patch grid, so changing the input size means resampling
them, and going straight from the pretraining resolution to the target resolution
would be too large a jump for that resampling. The first stage therefore trains at an
intermediate resolution, and the second stage begins from its best checkpoint and
trains at the full target resolution, resampling the position embeddings again.

One group of runs replaces the objective entirely. Instead of cross entropy they
optimize an AUC margin loss \cite{yuan2021dam}, which approximates the headline metric
directly, using the min-max optimizer that this loss requires. Mixed precision is disabled there and
the model is warm started from a cross entropy model, both following the practice
recommended for that loss. This changes the loss, the optimizer and the learning rate
scale together, so it is not a single variable comparison and is reported as such.

\subsection{Evaluation}

The headline metric is the mean AUROC over the five pathologies, with the per class
values and AUPRC also reported. Scoring is at the image level.

F1, precision, recall and specificity use one threshold per pathology chosen to
maximize F1 on val19k. Those thresholds are then frozen and applied unchanged to
val200 and test500.

During training the model is validated every few hundred optimizer steps on both
val19k and val200. The checkpoint that is kept is the one with the best mean AUROC on
the monitored set, and early stopping watches the same quantity.

The final ensampling results carry bootstrap confidence intervals over 10,000 resamples
of the images.

\section{Results}

\subsection{Architecture}
\label{sec:results-arch}

Table \ref{tab:arch} compares the backbones under the shared baseline setting, where
everything except the architecture is held fixed.

On val19k the backbones span 0.7994 to 0.8064, a range of 0.0070 across architectures
that differ by a factor of twenty eight in parameter count. On the radiologist sets the
same models span 0.0245 on val200 and 0.0170 on test500, around three times the seed
spread most configurations show in Section \ref{sec:results-noise}. The distance between
the best and the weakest backbone here is therefore a real difference, even though
neighbouring rows are not separable.

Nothing about capacity orders the results. ConvNeXt-L, at 196.2 million parameters,
finishes below ConvNeXt-T at 27.8 million on both radiologist sets, and DenseNet-121 at
7.0 million lands within 0.0042 of ConvNeXt-L on test500. ResNet-50, the oldest design
in the table, is fourth on test500 and ahead of both ConvNeXt-B and ConvNeXt-L.

The sets do not agree on a winner. ConvNeXt-L is first on val19k, while ConvNeXt-T leads
both radiologist sets, at 0.8850 and 0.8896. The two Swin rows sit at the bottom of
test500 despite being among the larger models. Most neighbouring rows differ by 0.002 to
0.005, which is too little to rank them, so only the distance between the top and the
bottom of the table should be read as a difference.

% Values from each run's results/<set>_frontal_results.json and training_summary.json.

\begin{table}[H]
  \caption{Backbones under the shared baseline setting. Resolution is
  $384\times320$ for every row except Swin-B, whose pretrained window size
  requires $384\times384$.}
  \centering
  \begin{tabular}{llcrrrr}
    \toprule
    Backbone & Pretraining & Pretrain res. & Params (M) & val19k & val200 & test500 \\
    \midrule
    ResNet-50 & IN-1k & $224^2$ & 23.5 & 0.7994 & 0.8759 & 0.8836 \\
    DenseNet-121 & IN-1k & $224^2$ & 7.0 & 0.8028 & 0.8727 & 0.8750 \\
    DenseNet-201 & IN-1k & $224^2$ & 18.1 & 0.8029 & 0.8755 & 0.8750 \\
    EfficientNetV2-M & IN-21k & $384^2$ & 52.9 & 0.8020 & 0.8766 & 0.8805 \\
    CoAtNet-1 & IN-12k & $224^2$ & 41.0 & 0.8026 & 0.8769 & 0.8824 \\
    Swin-B & IN-1k & $384^2$ & 86.9 & 0.8022 & 0.8605 & 0.8726 \\
    Swin-B & IN-22k & $384^2$ & 86.9 & 0.8036 & 0.8638 & 0.8761 \\
    ConvNeXt-T & IN-1k & $224^2$ & 27.8 & 0.8043 & 0.8850 & 0.8896 \\
    ConvNeXt-S & IN-1k & $224^2$ & 49.5 & 0.8056 & 0.8675 & 0.8789 \\
    ConvNeXt-B & IN-22k & $384^2$ & 87.6 & 0.8062 & 0.8806 & 0.8775 \\
    ConvNeXt-L & IN-22k & $384^2$ & 196.2 & 0.8064 & 0.8748 & 0.8792 \\
    \bottomrule
  \end{tabular}
  \label{tab:arch}
\end{table}


\subsection{Pretraining source}
\label{sec:results-pretrain}

Table \ref{tab:pretrain} holds the resolution fixed at the shared baseline and varies
only what the backbone saw before CheXpert.

The ImageNet label set makes no difference that survives the noise. Swin-B gains 0.0033
on val200 and 0.0035 on test500 from the larger label set, both a fraction of the seed
spread in Section \ref{sec:results-seed}. Pretraining on ten times more natural image
labels does not carry into this task.

An intermediate stage on ChestX-ray14 behaves differently depending on how that stage
is run. Copying the CheXpert schedule into it leaves ConvNeXt-B where it started, 0.8795
against 0.8775 on test500 and lower on val200. Running the same stage gently, at a tenth
of the learning rate for four epochs, gives 0.8942 on test500, 0.0167 above the baseline,
and softening both stages gives 0.8899. ConvNeXt-L moves the same way but by less,
gaining 0.0055 on val200 and 0.0019 on test500. Only the gentle ConvNeXt-B stage clears
the noise floor. One explanation is that the harm comes from the intermediate stage
overwriting the ImageNet features rather than from the extra data itself.

The backbones that were pretrained on chest radiographs from the start are the clearest
result in this section. All three Medical-MAE checkpoints beat every ImageNet backbone
in the table on test500, with the raw self-supervised checkpoint reaching 0.9046.
Against the matched ConvNeXt-B, which uses the same uncertainty policy, label smoothing
and resolution, the gain is 0.0061 on test500 and 0.0061 on val200. It costs nothing at
inference, since this ViT-B is 86.0 million parameters against 87.6 million for
ConvNeXt-B.

The three Medical-MAE checkpoints also show how little the ranking sets agree. The
checkpoint already fine tuned on CheXpert is far ahead on val19k at 0.8161, a margin of
0.0076 over the next best, and yet it is last of the three on both radiologist sets. The
raw checkpoint, which was never fine tuned on any labelled chest X-ray task, is the one
that transfers best. Section \ref{sec:results-valsize} returns to this disagreement.

% Values from each run's results/<set>_frontal_results.json; val19k from the run logs.

\begin{table}[H]
  \caption{Pretraining source at the shared $384\times320$ setting
  ($384\times384$ for Swin-B). Gentle stage 1 means the ChestX-ray14
  stage ran at $1\times10^{-5}$ for 4 epochs instead of $1\times10^{-4}$
  for 10, with the CheXpert stage left untouched. The last four rows use
  the mixed uncertainty policy and label smoothing, so the ConvNeXt-B row
  is repeated there under those same settings as the matched reference.}
  \centering
  \footnotesize
  \begin{tabular}{llccrrr}
    \toprule
    Backbone & Pretraining & U-policy & LS & val19k & val200 & test500 \\
    \midrule
    \multicolumn{7}{l}{\textit{ImageNet label set}} \\
    ConvNeXt-T & IN-1k & U-Ones & no & 0.8043 & 0.8850 & 0.8896 \\
    Swin-B & IN-1k & U-Ones & no & 0.8022 & 0.8605 & 0.8726 \\
    Swin-B & IN-22k & U-Ones & no & 0.8036 & 0.8638 & 0.8761 \\
    \midrule
    \multicolumn{7}{l}{\textit{ChestX-ray14 pretraining}} \\
    ConvNeXt-B & IN-22k & U-Ones & no & 0.8062 & 0.8806 & 0.8775 \\
    ConvNeXt-B & IN-22k + CXR14 & U-Ones & no & 0.8061 & 0.8759 & 0.8795 \\
    ConvNeXt-B & IN-22k + CXR14, gentle stage 1 & U-Ones & no & 0.8062 & 0.8919 & 0.8942 \\
    ConvNeXt-B & IN-22k + CXR14, gentle both stages & U-Ones & no & 0.8070 & 0.8895 & 0.8899 \\
    ConvNeXt-L & IN-22k & U-Ones & no & 0.8064 & 0.8748 & 0.8792 \\
    ConvNeXt-L & IN-22k + CXR14 & U-Ones & no & 0.8067 & 0.8803 & 0.8811 \\
    \midrule
    \multicolumn{7}{l}{\textit{Chest X-ray pretrained backbones}} \\
    ViT-B & Medical-MAE & mixed & yes & 0.8065 & 0.9030 & 0.9046 \\
    ViT-B & Medical-MAE + CXR14 & mixed & yes & 0.8085 & 0.8970 & 0.9020 \\
    ViT-B & Medical-MAE + CheXpert & mixed & yes & 0.8161 & 0.8906 & 0.8994 \\
    ConvNeXt-B & IN-22k & mixed & yes & 0.8038 & 0.8969 & 0.8985 \\
    \bottomrule
  \end{tabular}
  \label{tab:pretrain}
\end{table}


\subsection{Resolution}
\label{sec:results-res}

Table \ref{tab:res} raises the input resolution while holding the rest of the recipe
fixed. This is the largest single lever in the study, and it is also the only one whose
direction the radiologist sets agree on.

Going from $384\times320$ to $768\times640$ gains ConvNeXt-B 0.0052 on test500 and
0.0025 on val200. ConvNeXt-L gains 0.0195 and 0.0261 over the same step, although that
pair also changes the uncertainty policy and label smoothing and so overstates the
effect. Medical-MAE gains 0.0047 on test500 and 0.0094 on val200.

Above roughly $768\times640$ the gains stop. ConvNeXt-B at $1600\times1312$ scores the
same 0.9037 on test500 as its own $768\times640$ run, and RAD-DINO loses 0.0054 going
to $1064\times896$. Only val200 keeps rising, which is a property of that set rather
than of the models, as Section \ref{sec:results-noise} shows.

% Values from each run's results/<set>_frontal_results.json.

\begin{table}[H]
  \caption{Input resolution. Within each block only the resolution changes, except
  the ConvNeXt-L pair, where the larger run also uses the mixed uncertainty policy
  and label smoothing. Time is wall clock on the GPU listed in
  Table \ref{tab:allruns}.}
  \centering
  \footnotesize
  \begin{tabular}{llrrrr}
    \toprule
    Backbone & Resolution & Time & val19k & val200 & test500 \\
    \midrule
    ConvNeXt-B & $384\times320$ & 1h\,31m & 0.8038 & 0.8969 & 0.8985 \\
    ConvNeXt-B & $768\times640$ & 4h\,16m & 0.8065 & 0.8994 & 0.9037 \\
    ConvNeXt-B & $1600\times1312$ & 9h\,46m & 0.8030 & 0.9063 & 0.9037 \\
    \midrule
    ConvNeXt-L & $384\times320$ & 2h\,25m & 0.8064 & 0.8748 & 0.8792 \\
    ConvNeXt-L & $768\times640$ & 3h\,18m & 0.8062 & 0.9009 & 0.8987 \\
    \midrule
    ViT-B, Medical-MAE & $384\times320$ & 1h\,11m & 0.8085 & 0.8970 & 0.9020 \\
    ViT-B, Medical-MAE & $768\times640$ & 3h\,12m & 0.8076 & 0.9064 & 0.9067 \\
    \midrule
    ViT-B, RAD-DINO & $784\times644$ & 5h\,15m & 0.8043 & 0.9016 & 0.9100 \\
    ViT-B, RAD-DINO & $1064\times896$ & 6h\,14m & 0.8021 & 0.9000 & 0.9046 \\
    \bottomrule
  \end{tabular}
  \label{tab:res}
\end{table}


\subsection{Training choices}
\label{sec:results-uncertain}

Table \ref{tab:train} covers what to do with the uncertain labels, what to optimize,
and whether to smooth the targets. These runs sit on three different backbones, so
only comparisons inside a block mean anything.

On ResNet-50, U-Zeros loses 0.0171 on val200 and 0.0040 on test500, so mapping uncertain
cells to negative is the weaker choice there. On DenseNet-121 both alternatives to plain
U-Ones help: sending Consolidation's uncertain cells to negative gains 0.0176 on test500,
and the full per task scheme gains 0.0228. Reading uncertain cells as positive, the
simplest policy of all, is not the best on either backbone but is never far off.

On ConvNeXt-B the picture is clearer. Every variant that stops training Consolidation
toward the labeler's convention beats the U-Ones baseline on test500: masking uncertain
cells out of the loss gives 0.8963, self-training all five tasks gives 0.8970, and
self-training Consolidation alone gives 0.9002, against 0.8775 for the baseline. That
last gain of 0.0227 is the largest single effect in this table. The curriculum variant,
which shows the uncertain rows first, gives most of it back at 0.8848.

Label smoothing of 0.05 changes almost nothing. Against the run it was added to, it
moves test500 by $-0.0017$ and val200 by $-0.0002$, both well inside the noise. It is
kept in the later runs.

Replacing cross entropy with the AUC margin loss also fails to pay. Warm started from a
cross entropy model it reaches 0.8836 on test500, above the U-Ones baseline it warm
started from but below several plain cross entropy runs. When it is trained alone,
without the cross entropy stage, it falls to 0.7548 on val19k and 0.8539 on test500, the
worst result in the study. Optimizing the metric directly is not a shortcut to it.

% Values from each run's results/<set>_frontal_results.json; val19k from the run logs.

\begin{table}[H]
  \caption{Uncertainty policy, objective and label smoothing. Only comparisons inside
  one backbone block are meaningful. Mixed means the per task scheme of Section
  \ref{sec:uncertain}, with three way heads on Cardiomegaly and Pleural Effusion.
  val19k is scored under each run's own policy, so it is not comparable across rows;
  the radiologist sets have no uncertain labels and are.}
  \centering
  \footnotesize
  \begin{tabular}{llrrr}
    \toprule
    Run & Uncertainty and objective & val19k & val200 & test500 \\
    \midrule
    \multicolumn{5}{l}{\textit{ResNet-50}} \\
    baseline & U-Ones & 0.7994 & 0.8759 & 0.8836 \\
    u\_zeros & U-Zeros & 0.7990 & 0.8588 & 0.8796 \\
    \midrule
    \multicolumn{5}{l}{\textit{DenseNet-121}} \\
    baseline & U-Ones & 0.8028 & 0.8727 & 0.8750 \\
    cons\_zeros & U-Ones, U-Zeros on Consolidation & 0.8085 & 0.8834 & 0.8926 \\
    u\_mixed & mixed, U-Zeros on Consolidation & 0.8159 & 0.8820 & 0.8978 \\
    \midrule
    \multicolumn{5}{l}{\textit{ConvNeXt-B}} \\
    baseline & U-Ones & 0.8062 & 0.8806 & 0.8775 \\
    certain\_only & uncertain masked out of the loss & 0.7983 & 0.8808 & 0.8963 \\
    u\_selftrain\_all & U-SelfTrained on all five tasks & 0.7983 & 0.8856 & 0.8970 \\
    u\_selftrain\_consol\_only & U-Ones, U-SelfTrained on Consolidation & 0.8038 & 0.8971 & 0.9002 \\
    u\_selftrain & mixed, U-SelfTrained on Consolidation & 0.8220 & 0.8977 & 0.8977 \\
    curriculum & as consol\_only, uncertain rows first & 0.7984 & 0.8812 & 0.8848 \\
    final\_stage1 & as consol\_only, plus label smoothing & 0.8038 & 0.8969 & 0.8985 \\
    aucm & U-Ones, AUC-M after cross entropy & 0.8069 & 0.8825 & 0.8836 \\
    aucm\_no\_ce & mixed, AUC-M alone & 0.7548 & 0.8471 & 0.8539 \\
    \bottomrule
  \end{tabular}
  \label{tab:train}
\end{table}


\subsection{How much of this is noise}
\label{sec:results-noise}
\label{sec:results-seed}
\label{sec:results-valsize}

Every number quoted so far should be read against the size of the differences that
appear when nothing changes at all. Table \ref{tab:seeds} rebuilds the same
configuration under different random seeds.

The noise floor depends on the set. On val19k, with 19,100 images, three ConvNeXt-B
replicates span 0.0010 and three DenseNet-121 replicates span 0.0005. On the
radiologist sets the same replicates span 0.0053 to 0.0156 on val200 and 0.0036 to
0.0170 on test500. So the DenseNet-121 block shows that a single unlucky seed can
move test500 by 0.0170 on its own.

The three sets also measure different things, which is why they keep disagreeing.
val19k is large and therefore stable, but its labels come from the automatic labeler
and its uncertain cells are always scored as positive, so it rewards models that
reproduce the labeler. This shows up directly: the Medical-MAE checkpoint already fine
tuned on CheXpert leads val19k by 0.0076 and is mid table on both radiologist sets,
while the run that stopped training toward the labeler convention altogether, the one
that masks uncertain cells out of the loss, gives up 0.0079 on val19k and gains 0.0188
on test500.

val200 and test500 carry radiologist consensus labels, which is what the task actually
asks for, but at 202 and 518 images they are too small to rank models that sit within
0.015 of each other. Their own disagreement is the proof: the ordering of the twelve
backbones does not survive the move from one set to the other. A useful thought for the
rest of this report is to read differences below roughly 0.015 on test500 as ties, and
to take a lead on any single set as a hint rather than a settled ranking.

% Values from each run's results/<set>_frontal_results.json.

\begin{table}[H]
  \caption{Seed replicates. Every run inside a block is identical except for the
  random seed. Spread is the largest value minus the smallest.}
  \centering
  \footnotesize
  \begin{tabular}{llrrr}
    \toprule
    Run & Seed & val19k & val200 & test500 \\
    \midrule
    \multicolumn{5}{l}{\textit{ConvNeXt-B, $384\times320$}} \\
    final\_stage1 & 42 & 0.8038 & 0.8969 & 0.8985 \\
    seed7 & 7 & 0.8034 & 0.8916 & 0.8949 \\
    seed1337 & 1337 & 0.8028 & 0.8924 & 0.8967 \\
    \multicolumn{2}{l}{\textit{spread}} & 0.0010 & 0.0053 & 0.0036 \\
    \midrule
    \multicolumn{5}{l}{\textit{DenseNet-121, $384\times320$}} \\
    densenet121 & 42 & 0.8028 & 0.8727 & 0.8750 \\
    seed7 & 7 & 0.8023 & 0.8740 & 0.8804 \\
    seed123 & 123 & -- & 0.8584 & 0.8634 \\
    \multicolumn{2}{l}{\textit{spread}} & 0.0005 & 0.0156 & 0.0170 \\
    \midrule
    \multicolumn{5}{l}{\textit{ViT-B Medical-MAE, $768\times640$}} \\
    nih\_B\_768\_s2 & 42 & 0.8076 & 0.9064 & 0.9067 \\
    seed7 & 7 & 0.8090 & 0.9045 & 0.9077 \\
    seed1337 & 1337 & 0.8053 & 0.9105 & 0.9112 \\
    \multicolumn{2}{l}{\textit{spread}} & 0.0037 & 0.0059 & 0.0046 \\
    \bottomrule
  \end{tabular}
  \label{tab:seeds}
\end{table}


\subsection{Ensembling and final result}
\label{sec:results-ens}

Table \ref{tab:ens} collects the ensembles. What decides their quality is not how many
members they hold but how different those members are.

Averaging five checkpoints of one ConvNeXt-B reaches 0.9076, and widening that to seven
ConvNeXt-B runs, including separate seeds, reaches 0.9095. Both sit below the best single
model, which scores 0.9113 on its own. Replacing those seven near copies with three
genuinely different backbones, ConvNeXt-B at $1600\times1312$, Medical-MAE at
$768\times640$ and RAD-DINO at $784\times644$, lifts a plain probability average to
\textbf{0.9174}, the best result in the study. Fitting per class weights on val200 does
not improve on that plain average.

Over 10,000 paired resamples of the 518 test500 images, that ensemble scores 0.9174 with
a 95\% interval of $[0.9040, 0.9303]$, and its gain over the strongest single member is
$+0.0061$ with a 95\% interval of $[+0.0005, +0.0118]$. Per class the ensemble
reaches 0.9712 on Pleural Effusion, 0.9269 on Edema, 0.9138 on Cardiomegaly, 0.9127 on
Consolidation and 0.8623 on Atelectasis.

% Values from training_scripts/others/ensembling_results/*/ensemble_test500_frontal_summary.json
% and training_scripts/others/ci/results/2026-08-15_14-24-38_test500_frontal/.

\begin{table}[H]
  \caption{Ensembles on test500. Weights, when used, are per class and fitted on
  val200. The single model row is the strongest member on its own.}
  \centering
  \footnotesize
  \begin{tabular}{llcr}
    \toprule
    Members & Blend & Weights & test500 \\
    \midrule
    Medical-MAE ViT-B $768\times640$ alone & -- & -- & 0.9113 \\
    \midrule
    5 ConvNeXt-B checkpoints & probability average & no & 0.9076 \\
    7 ConvNeXt-B runs & probability average & no & 0.9095 \\
    \midrule
    3 backbones & probability average & no & \textbf{0.9174} \\
    8 backbones & logit average & yes & 0.9159 \\
    \bottomrule
  \end{tabular}
  \label{tab:ens}
\end{table}


\section{Discussion}

Across roughly fifty runs, a few choices moved the metric and most of the rest did not.

Resolution moved it, up to a point. Going from $384\times320$ to $768\times640$ was the
largest single gain in the study, and going further gave almost nothing back. Domain
pretraining moved it: backbones pretrained on chest radiographs beat
every ImageNet backbone at equal size and equal resolution, which is the one place where
the choice of starting weights clearly paid. Ensembling moved it, but only when the
members were genuinely different from one another.

Architecture moved it too, but not in the way parameter counts suggest. The eleven
backbones span 0.0170 on test500 and 0.0245 on val200, around three times the seed
spread that most configurations show, so the difference between a good backbone and a
weak one is real. What does not follow is capacity: ConvNeXt-T at 27.8 million parameters finishes above
ConvNeXt-L at 196.2 million on both radiologist sets, and on test500 even ResNet-50
lands ahead of ConvNeXt-B and ConvNeXt-L. Picking a good family matters, and paying for
a larger member of that family does not.

What did not move the metric is a shorter list than expected. Enlarging the ImageNet
label set from 1k to 22k changed nothing that survives the noise, on either    backbone
where both were available. That is a statement about the label set alone: every
backbone here starts from pretrained weights and none was trained from scratch, so this
study says nothing about the value of ImageNet pretraining itself, which is almost
certainly large. The uncertainty policy, which the CheXpert literature has spent
considerable effort on, mattered much less than the seed. Label smoothing did nothing
measurable. Optimizing AUROC directly through an AUC margin loss did not beat cross
entropy, and without a cross entropy stage it collapsed.

\section{Limitations}

Everything here is one dataset. There is no external validation on MIMIC-CXR,
ChestX-ray14 as a target, or any prospective set, so none of these conclusions are known
to transfer.

The final ensemble is selected and weighted on val200, which holds 202 images and, for
Consolidation, 32 positives. The bootstrap interval on test500 accounts for sampling the
test images but not for that selection, so any headline figure on test500 could be
optimistic, and this applies to any published number on that set. The per class weights should be read as one draw from a noisy
procedure rather than as a finding about which model is good at which pathology.

Not every ablation shares one backbone. The uncertainty policies are spread across
ResNet-50, DenseNet-121 and ConvNeXt-B, and the ChestX-ray14 pretraining runs vary the
schedule alongside the extra stage, so a few comparisons in Section
\ref{sec:results-uncertain} and Section \ref{sec:results-pretrain} carry more than one
variable. Only three configurations were repeated across seeds, so the noise floor is
estimated from a small sample and may be understated for the configurations that were
never repeated.

\section{Conclusion}

Under one fixed pipeline on CheXpert, input resolution, chest X-ray pretraining and the
handling of uncertain labels were the design choices that paid, ensembling paid when its
members were diverse, and the architecture mattered while its size did not as much. The
remaining choices, including the ImageNet label set, label smoothing and the objective,
were within the noise of rerunning the same configuration with a different seed. The best
single model reached 0.9113 mean AUROC on the official test set and a plain average of
three different backbones reached 0.9174, with a paired bootstrap interval of
$[+0.0005, +0.0118]$ on the difference. The wider point is that the official radiologist
sets are small enough that model rankings drawn from them are unstable, and reported
gains below roughly 0.015 should be treated as ties.

\bibliographystyle{unsrt}
\bibliography{references}

\appendix
\section{Run configurations}
\label{app:runs}

Epochs are the configured budget, with the total number of optimizer steps in
parentheses. Training stops early, so the selected checkpoint is usually well
before the end. Batch size is 64 everywhere except
\texttt{\scriptsize convnext\_base\_22k\_1600x1312}, which used 32.

\begin{landscape}
  \scriptsize
  \setlength{\tabcolsep}{1.5pt}
  \begin{longtable}{llllllllllrrr}
    \caption{Configuration and results of every run reported in this study.}
    \label{tab:allruns} \\
    \toprule
    Run & Pretraining & Resolution & LR peak & LR min & Epochs (steps) & GPU & U-policy & LS & Time & val19k & val200 & test500 \\
    \midrule
    \endfirsthead
    \multicolumn{13}{l}{\textit{Table \thetable\ continued from the previous page}} \\
    \toprule
    Run & Pretraining & Resolution & LR peak & LR min & Epochs (steps) & GPU & U-policy & LS & Time & val19k & val200 & test500 \\
    \midrule
    \endhead
    \midrule
    \multicolumn{13}{r}{\textit{continued on the next page}} \\
    \endfoot
    \bottomrule
    \endlastfoot
    \midrule
    \multicolumn{13}{l}{\textit{Architecture}} \\
    \texttt{\scriptsize resnet50\_without\_clahe} & IN-1k & 384$\times$320 & $1\times10^{-4}$ & $1\times10^{-6}$ & 15 (40,305) & L4 & U-Ones & no & 2h\,08m & 0.7994 & 0.8759 & 0.8836 \\
    \texttt{\scriptsize densenet121} & IN-1k & 384$\times$320 & $1\times10^{-4}$ & $1\times10^{-6}$ & 10 (26,870) & L4 & U-Ones & no & 2h\,15m & 0.8028 & 0.8727 & 0.8750 \\
    \texttt{\scriptsize densenet201} & IN-1k & 384$\times$320 & $1\times10^{-4}$ & $1\times10^{-6}$ & 10 (26,870) & L4 & U-Ones & no & 3h\,06m & 0.8029 & 0.8755 & 0.8750 \\
    \texttt{\scriptsize convnext\_tiny} & IN-1k & 384$\times$320 & $1\times10^{-4}$ & $1\times10^{-6}$ & 10 (26,870) & L4 & U-Ones & no & 1h\,58m & 0.8043 & 0.8850 & 0.8896 \\
    \texttt{\scriptsize convnext\_tiny\_22k} & IN-22k & 384$\times$320 & $1\times10^{-4}$ & $1\times10^{-6}$ & 10 (26,870) & L4 & U-Ones & no & 2h\,07m & 0.8045 & 0.8788 & 0.8835 \\
    \texttt{\scriptsize convnext\_small} & IN-1k & 384$\times$320 & $1\times10^{-4}$ & $1\times10^{-6}$ & 10 (26,870) & L4 & U-Ones & no & 3h\,17m & 0.8056 & 0.8675 & 0.8789 \\
    \texttt{\scriptsize convnext\_base\_22k} & IN-22k & 384$\times$320 & $1\times10^{-4}$ & $1\times10^{-6}$ & 10 (26,870) & A100 & U-Ones & no & 1h\,17m & 0.8062 & 0.8806 & 0.8775 \\
    \texttt{\scriptsize convnext\_large\_22k} & IN-22k & 384$\times$320 & $1\times10^{-4}$ & $1\times10^{-6}$ & 10 (26,870) & A100 & U-Ones & no & 2h\,25m & 0.8064 & 0.8748 & 0.8792 \\
    \texttt{\scriptsize swin\_base} & IN-1k & 384$\times$384 & $1\times10^{-4}$ & $1\times10^{-6}$ & 10 (26,870) & A100-80GB & U-Ones & no & 2h\,17m & 0.8022 & 0.8605 & 0.8726 \\
    \texttt{\scriptsize swin\_base\_22k} & IN-22k & 384$\times$384 & $1\times10^{-4}$ & $1\times10^{-6}$ & 10 (26,870) & A100-80GB & U-Ones & no & 2h\,29m & 0.8036 & 0.8638 & 0.8761 \\
    \texttt{\scriptsize efficientnetv2\_m} & IN-21k & 384$\times$320 & $1\times10^{-4}$ & $1\times10^{-6}$ & 10 (26,870) & A100 & U-Ones & yes & 1h\,18m & 0.8020 & 0.8766 & 0.8805 \\
    \texttt{\scriptsize coatnet\_1\_12k} & IN-12k & 384$\times$320 & $1\times10^{-4}$ & $1\times10^{-6}$ & 10 (26,870) & A100 & mixed & yes & 1h\,54m & 0.8026 & 0.8769 & 0.8824 \\
    \midrule
    \multicolumn{13}{l}{\textit{Preprocessing}} \\
    \texttt{\scriptsize resnet50\_with\_clahe} & IN-1k & 384$\times$320 & $1\times10^{-4}$ & $1\times10^{-6}$ & 15 (40,305) & L4 & U-Ones & no & 3h\,39m & 0.7973 & 0.8664 & 0.8732 \\
    \midrule
    \multicolumn{13}{l}{\textit{Uncertain labels}} \\
    \texttt{\scriptsize resnet50\_u\_zeros} & IN-1k & 384$\times$320 & $1\times10^{-4}$ & $1\times10^{-6}$ & 15 (40,305) & L4 & U-Zeros & no & 2h\,38m & 0.7990 & 0.8588 & 0.8796 \\
    \texttt{\scriptsize densenet121\_cons\_zeros} & IN-1k & 384$\times$320 & $1\times10^{-4}$ & $1\times10^{-6}$ & 10 (26,870) & L4 & U-Ones & no & 2h\,13m & 0.8085 & 0.8834 & 0.8926 \\
    \texttt{\scriptsize densenet121\_u\_mixed} & IN-1k & 384$\times$320 & $1\times10^{-4}$ & $1\times10^{-6}$ & 10 (26,870) & L4 & mixed & no & 2h\,39m & 0.8159 & 0.8820 & 0.8978 \\
    \texttt{\scriptsize convnext\_base\_22k\_certain\_only} & IN-22k & 384$\times$320 & $1\times10^{-4}$ & $1\times10^{-6}$ & 10 (26,870) & A100 & U-Ones & no & 1h\,28m & 0.7983 & 0.8808 & 0.8963 \\
    \texttt{\scriptsize convnext\_base\_22k\_u\_selftrain} & IN-22k & 384$\times$320 & $1\times10^{-4}$ & $1\times10^{-6}$ & 10 (26,870) & A100 & mixed & no & 1h\,47m & 0.8220 & 0.8977 & 0.8977 \\
    \texttt{\scriptsize convnext\_base\_22k\_u\_selftrain\_all} & IN-22k & 384$\times$320 & $1\times10^{-4}$ & $1\times10^{-6}$ & 10 (26,870) & L4 & mixed & no & 3h\,37m & 0.7983 & 0.8856 & 0.8970 \\
    \texttt{\scriptsize convnext\_base\_22k\_u\_selftrain\_consol\_only} & IN-22k & 384$\times$320 & $1\times10^{-4}$ & $1\times10^{-6}$ & 10 (26,870) & A100 & mixed & no & 1h\,21m & 0.8038 & 0.8971 & 0.9002 \\
    \texttt{\scriptsize convnext\_base\_22k\_curriculum} & IN-22k & 384$\times$320 & $1\times10^{-4}$ & $1\times10^{-6}$ & 10 (16,818) & A100 & mixed & no & 58m\,11s & 0.7984 & 0.8812 & 0.8848 \\
    \midrule
    \multicolumn{13}{l}{\textit{Label smoothing}} \\
    \texttt{\scriptsize convnext\_base\_22k\_final\_stage1} & IN-22k & 384$\times$320 & $1\times10^{-4}$ & $1\times10^{-6}$ & 10 (26,870) & A100 & mixed & yes & 1h\,31m & 0.8038 & 0.8969 & 0.8985 \\
    \midrule
    \multicolumn{13}{l}{\textit{Objective}} \\
    \texttt{\scriptsize convnext\_base\_22k\_aucm} & IN-22k & 384$\times$320 & $0.075$ & $7.5\times10^{-4}$ & 6 (16,122) & A100 & U-Ones & no & 1h\,15m & 0.8069 & 0.8825 & 0.8836 \\
    \texttt{\scriptsize convnext\_base\_22k\_aucm\_no\_ce} & IN-22k & 384$\times$320 & $0.075$ & $7.5\times10^{-4}$ & 6 (16,122) & A100 & mixed & yes & 1h\,40m & 0.7548 & 0.8471 & 0.8539 \\
    \midrule
    \multicolumn{13}{l}{\textit{Domain pretraining}} \\
    \texttt{\scriptsize medmae\_vitb\_raw} & MAE & 384$\times$320 & $5\times10^{-5}$ & $5\times10^{-7}$ & 6 (16,122) & A100 & mixed & yes & 1h\,03m & 0.8065 & 0.9030 & 0.9046 \\
    \texttt{\scriptsize medmae\_vitb\_nih} & MAE+CXR14 & 384$\times$320 & $5\times10^{-5}$ & $5\times10^{-7}$ & 6 (16,122) & A100 & mixed & yes & 1h\,11m & 0.8085 & 0.8970 & 0.9020 \\
    \texttt{\scriptsize medmae\_vitb\_chexpert} & MAE+CheXpert & 384$\times$320 & $1\times10^{-5}$ & $1\times10^{-7}$ & 6 (16,122) & A100 & mixed & yes & 1h\,04m & 0.8161 & 0.8906 & 0.8994 \\
    \texttt{\scriptsize convnext\_base\_22k\_cxr14\_pretrain} & IN-22k+CXR14 & 384$\times$320 & $1\times10^{-4}$ & $1\times10^{-6}$ & 10 (26,870) & A100 & U-Ones & no & 1h\,33m & 0.8061 & 0.8759 & 0.8795 \\
    \texttt{\scriptsize convnext\_base\_22k\_cxr14\_pretrain\_lowlr} & IN-22k+CXR14 & 384$\times$320 & $1\times10^{-5}$ & $1\times10^{-7}$ & 4 (26,870) & A100 & U-Ones & no & 1h\,12m & 0.8062 & 0.8919 & 0.8942 \\
    \texttt{\scriptsize convnext\_base\_22k\_cxr14\_pretrain\_lowlr\_all} & IN-22k+CXR14 & 384$\times$320 & $1\times10^{-5}$ & $1\times10^{-7}$ & 4 (16,122) & A100 & U-Ones & no & 1h\,35m & 0.8070 & 0.8895 & 0.8899 \\
    \texttt{\scriptsize convnext\_large\_22k\_cxr14\_pretrain} & IN-22k+CXR14 & 384$\times$320 & $1\times10^{-4}$ & $1\times10^{-6}$ & 10 (26,870) & A100 & U-Ones & no & 1h\,30m & 0.8067 & 0.8803 & 0.8811 \\
    \midrule
    \multicolumn{13}{l}{\textit{Resolution}} \\
    \texttt{\scriptsize convnext\_base\_22k\_768x640} & IN-22k & 768$\times$640 & $1\times10^{-4}$ & $1\times10^{-6}$ & 10 (26,870) & A100-80GB & mixed & yes & 4h\,16m & 0.8065 & 0.8994 & 0.9037 \\
    \texttt{\scriptsize convnext\_base\_22k\_1600x1312} & IN-22k & 1600$\times$1312 & $1\times10^{-4}$ & $1\times10^{-6}$ & 10 (26,870) & B200 & mixed & yes & 9h\,46m & 0.8030 & 0.9063 & 0.9037 \\
    \texttt{\scriptsize convnext\_large\_22k\_768x640} & IN-22k & 768$\times$640 & $1\times10^{-4}$ & $1\times10^{-6}$ & 10 (26,870) & H200 & mixed & yes & 3h\,18m & 0.8062 & 0.9009 & 0.8987 \\
    \texttt{\scriptsize medmae\_vitb\_nih\_B\_768\_s2} & MAE+CXR14 & 768$\times$640 & $5\times10^{-5}$ & $5\times10^{-7}$ & 8 (21,496) & A100-80GB & mixed & yes & 3h\,12m & 0.8076 & 0.9064 & 0.9067 \\
    \texttt{\scriptsize rad\_dino\_vitB\_768} & DINOv2 & 784$\times$644 & $3\times10^{-5}$ & $3\times10^{-7}$ & 5 (13,435) & A100-80GB & mixed & yes & 5h\,15m & 0.8044 & 0.9016 & 0.9100 \\
    \texttt{\scriptsize rad\_dino\_vitB\_1064x896} & DINOv2 & 1064$\times$896 & $6\times10^{-6}$ & $1\times10^{-7}$ & 6 (16,122) & H200 & mixed & yes & 6h\,14m & 0.8021 & 0.9000 & 0.9046 \\
    \midrule
    \multicolumn{13}{l}{\textit{Seed replicates}} \\
    \texttt{\scriptsize convnext\_base\_22k\_seed7} & IN-22k & 384$\times$320 & $1\times10^{-4}$ & $1\times10^{-6}$ & 10 (26,870) & A100 & mixed & yes & 51m\,18s & 0.8034 & 0.8916 & 0.8949 \\
    \texttt{\scriptsize convnext\_base\_22k\_seed1337} & IN-22k & 384$\times$320 & $1\times10^{-4}$ & $1\times10^{-6}$ & 10 (26,870) & A100 & mixed & yes & 54m\,20s & 0.8028 & 0.8924 & 0.8967 \\
    \texttt{\scriptsize densenet121\_seed7} & IN-1k & 384$\times$320 & $1\times10^{-4}$ & $1\times10^{-6}$ & 10 (26,870) & L4 & U-Ones & no & 2h\,28m & 0.8023 & 0.8740 & 0.8804 \\
    \texttt{\scriptsize densenet121\_seed123} & IN-1k & 384$\times$320 & $1\times10^{-4}$ & $1\times10^{-6}$ & 10 (26,870) & L4 & U-Ones & no & 2h\,41m & 0.8023 & 0.8584 & 0.8634 \\
    \texttt{\scriptsize medmae\_vitb\_nih\_B\_448\_s1\_seed7} & MAE+CXR14 & 448$\times$384 & $5\times10^{-5}$ & $1\times10^{-5}$ & 1 (2,687) & A100 & mixed & yes & 39m\,24s & 0.7960 & 0.8972 & 0.9080 \\
    \texttt{\scriptsize medmae\_vitb\_nih\_B\_448\_s1\_seed1337} & MAE+CXR14 & 448$\times$384 & $5\times10^{-5}$ & $1\times10^{-5}$ & 1 (2,687) & A100 & mixed & yes & 33m\,06s & 0.8019 & 0.9052 & 0.9088 \\
    \texttt{\scriptsize medmae\_vitb\_nih\_B\_768\_s2\_seed7} & MAE+CXR14 & 768$\times$640 & $5\times10^{-5}$ & $5\times10^{-7}$ & 8 (21,496) & A100-80GB & mixed & yes & 2h\,22m & 0.8090 & 0.9045 & 0.9077 \\
    \texttt{\scriptsize medmae\_vitb\_nih\_B\_768\_s2\_seed1337} & MAE+CXR14 & 768$\times$640 & $5\times10^{-5}$ & $5\times10^{-7}$ & 8 (21,496) & A100-80GB & mixed & yes & 2h\,06m & 0.8053 & 0.9105 & 0.9112 \\
  \end{longtable}
\end{landscape}

\section{Per-disease runs}
\label{app:perdisease}

These runs train a single pathology rather than all five, so their scores are
single class AUROC and are not comparable to the mean AUROC reported elsewhere.
The first group trains one ConvNeXt-B per pathology with cross entropy. The
second group replaces the objective with the AUC margin loss, warm started
from a shared cross entropy backbone.

\begin{table}[H]
  \caption{Single pathology runs. Scores are single class AUROC.}
  \centering
  \footnotesize
  \setlength{\tabcolsep}{4pt}
  \begin{tabular}{lllrrr}
    \toprule
    Run & Pathology & Objective & val19k & val200 & test500 \\
    \midrule
    \texttt{\scriptsize convnext\_base\_22k\_atelectasis\_only} & Atelectasis & BCE & 0.7243 & 0.8500 & 0.8431 \\
    \texttt{\scriptsize convnext\_base\_22k\_cardiomegaly\_only} & Cardiomegaly & BCE & 0.8716 & 0.8277 & 0.8850 \\
    \texttt{\scriptsize convnext\_base\_22k\_consolidation\_only} & Consolidation & BCE & 0.6800 & 0.8947 & 0.9155 \\
    \texttt{\scriptsize convnext\_base\_22k\_edema\_only} & Edema & BCE & 0.8517 & 0.9298 & 0.9153 \\
    \texttt{\scriptsize convnext\_base\_22k\_pleural\_effusion\_only} & Pleural Effusion & BCE & 0.8782 & 0.9299 & 0.9495 \\
    \texttt{\scriptsize convnext\_base\_22k\_final\_stage2\_atelectasis} & Atelectasis & AUC-M & 0.7233 & 0.8557 & 0.8363 \\
    \texttt{\scriptsize convnext\_base\_22k\_final\_stage2\_cardiomegaly} & Cardiomegaly & AUC-M & 0.8712 & 0.8319 & 0.8949 \\
    \texttt{\scriptsize convnext\_base\_22k\_final\_stage2\_consolidation} & Consolidation & AUC-M & 0.6820 & 0.9175 & 0.8851 \\
    \texttt{\scriptsize convnext\_base\_22k\_final\_stage2\_edema} & Edema & AUC-M & 0.8485 & 0.9454 & 0.9232 \\
    \texttt{\scriptsize convnext\_base\_22k\_final\_stage2\_pleural\_effusion} & Pleural Effusion & AUC-M & 0.8767 & 0.9324 & 0.9557 \\
    \bottomrule
  \end{tabular}
  \label{tab:perdisease}
\end{table}




\end{document}
